\section{Implications for behavioral evaluation}
\label{sec:discussion}

\subsection{What the eight questions establish here}

For strategic reasoning and economic risk (Q1 and Q6), the strongest answers concern tradeoffs and accessible resources. Recipients do not automatically accept ownership completion: the complete offer population contains 74 direct refusals of such offers, 73 without breaking another complete recipient set. Their recorded explanations include consequences for rivals as well as their own costs. Legal-sequence checks independently distinguish current cash, reachable finance and nominal assets. These findings establish selective behavior and particular financial successes or mistakes. They do not establish efficient prices, optimal risk tolerance or maximized winning probability.

For opponent modeling and sustained plans (Q2 and Q4), the temporal records contain prospective statements with different kinds of targets. Some concern observable financial capacity, some a rival's next action, and others an open-ended intention. A target must be specified and an opportunity must occur before an outcome is scored. The G9 development objective is eventually realized; the G16 Virginia purchase encounters a legal opportunity but the received action starts an auction; several explicit landing policies encounter no opportunity. These outcomes are informative about execution. They do not identify a private-belief model, calibrated theory of mind or the causal value of remembering an earlier statement.

For negotiation and collective dynamics (Q3 and Q7), the complete thread and transfer ledgers establish who exchanges what, how often counters occur, and how rent and inheritance concentrate resources. The transcripts add the sequence of appeals, corrections, omissions and withdrawals. They show documented coordination and misleading communication without turning dense trade links into collusion. A causal account of agreement, a coalition-stability measure or a welfare comparison needs additional definitions and interventions. The eventual winner's receipt of 202 of 227 nonterminal inherited deeds is a substantial realized concentration, but does not isolate an inheritance effect on victory.

For adaptation and validity (Q5 and Q8), we can observe revision and test whether a specific error recurs. The 4/44 recurrence after an applied construction repair differs from the 9/50 recurrence after a recorded corrective request; neither is a learning treatment effect. Internal validity is stronger where complete populations, replay and executable endpoints are available. External validity remains weak: the LLM cohort holds opponent identities fixed, and the separate baseline pilot reveals sparse development, no trade initiation and extensive censoring. The scenario audit further shows that existing preferred actions and explanation rubrics need adjudication before serving as an independent skill test.

These answers also separate the three broader questions. The study gives direct counterexamples to equating strategic-sounding text with strategic execution (A). It does not estimate the association between general and multi-agent capability (B): there is no independent capability axis or held-out family test. It identifies some routes by which mistakes and inheritance contribute to outcomes, but does not decompose winning into skill, chance and opponent weakness (C). Appendix~\ref{sec:agenda-prospective-evaluation} sets out distinct experiments for the unresolved parts. They remain unexecuted designs, rather than results implied by this audit.

\subsection{The label determines which evidence is missing}

The same END can be a valid selected response, a deterministic fallback, or a received call inconsistent with accompanying content. These distinctions do not remove failures from the configured agent's performance. A defaulted build still loses an opportunity in the observed game. They change the explanation: a failure to execute an attempted allocation is not evidence of preferring reserves. Evaluations should report system outcomes together with their resolution provenance, rather than either attributing every outcome to strategy or excluding implementation-mediated failures.

Opportunity conditioning also changes the question. The 43/63 first-window construction count retains prior acquisition and finance; the 34/42 entry-buildable count describes a narrower condition already reached. The nine groups never buildable during the first window and the sixty-nine terminal inherited completions are different kinds of missing opportunity. A single conversion rate hides both. Our per-game ledger retains zero-denominator cells and seed blocks so thousands of dependent decisions cannot be mistaken for thousands of independent trials.

Feasibility is another distinct label. A current menu can understate future collateral unlocked by sales, while nominal wealth can overstate access under the hotel rule. Positive and negative claims have asymmetric evidence requirements: one legal witness establishes a reachable endpoint, whereas one failed search establishes nothing. The 3/60 payable eliminations and the 1/60 nominally solvent hotel lock demonstrate both directions in this corpus. None establishes whether continuing was preferable or would change the winner.

These results validate specific classifications and refute specific sufficiency rules. They do not measure how much a general annotation system improves after adopting the protocol. That would require a frozen target population, prespecified labels and independently adjudicated evaluation. The selected social contrasts are particularly unsuited to an error-rate claim. Their role is to show concrete alternatives that a label must exclude, such as acceptance with explicit premise correction.

\subsection{Norms and recipient effects require separate designs}
\label{sec:discussion-design}

The motivating capability--alignment question needs more than four win totals. This cohort has four configured systems, unequal response conditions and no independently measured capability axis. Public deception is permitted; an external dishonesty taxonomy would therefore answer a different question from adherence to the assigned prompt. Private honesty also refers to reasoning that a transcript does not directly reveal. A false note can be mistaken, a transient statement can be corrected, and a received action can contradict its own explanation.

The evidence still supplies useful experimental cases. A disclosure task could fix a rule requiring revelation of a specified encumbrance and vary the incentive to omit it. A standing-offer task could fix terms, expiry and release authority before a G18-style invocation. Recipient-effect experiments could randomize a message or its financial framing at a verified historical state, repeat continuations, and distinguish correct factual restatement, action choice and economic outcome. Removing a message is one possible intervention, not a uniquely sufficient test of persuasion. These would be new versioned task conditions, not norms retroactively assigned to the existing games.

For comparisons across capability, an independently measured axis with several tiers within families would avoid identifying capability with the same outcomes used to assess conduct. Prompts, memory, retry policies, tool choice and resources should be fixed or deliberately varied. Economic utility, instruction adherence and recipient effects should remain separate outcomes. This is consistent with mixed existing findings: Arena reports no capability association under its economic ordering and codebook, while Emergent Collusion reports within-family differences under designed violation opportunities \citep{li2026arena,shi2026emergentcollusion}. MASK separates accuracy from honesty; Safetywashing questions whether purported safety measures mainly track capability \citep{ren2025mask,ren2024safetywashing}. The present cohort adds no point to a capability--alignment curve.

Exact-prefix certificates can support held-out tasks with an explicit payment or construction endpoint. A monitor study would additionally fix which channels are visible, the prediction target and scoring rule. Reinforcement learning would require validated reward semantics, observation/action adapters, reset distributions and horizon rules. Existing replay does not supply those experiments automatically. Both the earlier capped prefix continuation and the nineteen censored baseline-only pilot games caution against assuming termination or treating available reference policies as adequate strategic opponents.

\subsection{The audit must survive attacks on its own assumptions}

Original-source checks corrected several of our labels: G15's productive group was orange rather than terminally inherited green; G19's cumulative construction was \$2,000 rather than a \$1,200 step; a transaction's initiator need not author its accepted terms. These are measurement errors even when the raw event log is correct. Agreement between fluent reviewers would not settle them.

The new verifier itself required correction. A fresh reviewer found that engine-accepted bundled sales can escape an atomic search, so the original generic negative-search claim was unsound. We preserve the counterexample, restrict graph certificates to the supported mortgage-only condition, and test refusal on the unsupported case. All historical classifications survive that change, but the broader claim does not. This is why the artifact records intervention classes, proof assumptions and unresolved outcomes rather than presenting an engine call as a universal oracle.

The resulting contribution is an inspectable link between a stated behavioral claim and evidence adequate for that claim. Historical replay checks what happened. Attempt provenance checks how an action was resolved. Legal witnesses and bounds check a specified reachable endpoint. Recipient records constrain what social conclusions follow. Keeping these tests separate makes both positive results and remaining uncertainty reviewable.

